\documentclass[11pt]{article}
\usepackage[margin=1in]{geometry}
\usepackage{amsmath}
\usepackage{booktabs}

\title{Gradient Boosting on a Small Tabular Benchmark: The Wine Dataset}
\author{Paper2Proof canary set\thanks{A small paper written for the Paper2Proof evaluation set (eval set A).}}
\date{2026}

\begin{document}
\maketitle

\begin{abstract}
We evaluate gradient-boosted decision trees with default capacity on the scikit-learn wine
dataset, a 178-sample, three-class chemical analysis task. Boosting reaches a test accuracy of
0.969 averaged over three random splits.
\end{abstract}

\section{Introduction}
Tree ensembles are the default choice for small tabular problems. We report a reproducible
reference number for gradient boosting on the wine dataset.

\section{Method}
Gradient-boosted decision trees for multiclass classification (log-loss): 100 boosting stages,
learning rate 0.1, maximum tree depth 3, no subsampling, with the trees' random state set to the
seed of the split. Features are used as provided (no scaling).

\section{Experimental setup}
\paragraph{Data.} The wine dataset shipped with scikit-learn (\texttt{load\_wine}): 178 samples,
13 features, 3 classes.
\paragraph{Splits.} For each seed $s \in \{0, 1, 2\}$, a stratified split with 30\,\% test
samples (\texttt{train\_test\_split(test\_size=0.3, random\_state=s, stratify=y)}).
\paragraph{Metric.} Test accuracy, mean and sample standard deviation over the three seeds.
\paragraph{Software.} Python 3.12, scikit-learn 1.9.1, CPU only.

\section{Results}
Table~\ref{tab:wine} reports the result: gradient boosting reaches a test accuracy of
0.969 $\pm$ 0.021 on wine. The spread reflects the small test set (54 samples).

\begin{table}[t]
\centering
\caption{Test accuracy on wine (mean $\pm$ standard deviation over seeds 0, 1, 2).}
\label{tab:wine}
\begin{tabular}{lc}
\toprule
Model & Accuracy \\
\midrule
Gradient boosting (100 stages, depth 3) & $0.969 \pm 0.021$ \\
\bottomrule
\end{tabular}
\end{table}

\section{Conclusion}
Default gradient boosting is close to the ceiling on wine; differences between strong models on
this dataset are within the split-to-split noise.

\end{document}
